\section{加速度と位置}\label{節:加速度と位置}
\chapterref{章:導入}では,速度を短い時間ごとに足して位置の変化を求めた.部分積分を使うと,同じ変化を加速度から表すこともできる.

速度$v(t)$は連続微分可能とし,加速度を$a(t)=v'(t)$とする.\,$a(t)$は速度がどの程度の割合で変わるかを表す量である.時刻$0$から$T$までの位置の変化を$D$とすると,
\[
 D=\int_0^T v(t)\dd t.
\]
ここで$t$の代わりに$t-T$を微分された因子と見て,部分積分してみよう.
\begin{align*}
 D&=[(t-T)v(t)]_0^T-\int_0^T(t-T)v'(t)\dd t\\
 &=Tv(0)+\int_0^T(T-t)a(t)\dd t.
\end{align*}
最初の項は,最初の速度$v(0)$のまま動いた場合の変化である.残りの項は,その後に速度が変わった分を表す.時刻$t$付近で速度が$a(t)\Delta t$だけ変われば,その変化は残りの時間$T-t$にわたって位置へ影響する.そのため,\,$T-t$が掛かると読むこともできる.

例えば加速度が一定値$a$なら,
\[
 D=Tv(0)+a\int_0^T(T-t)\dd t
 =Tv(0)+\frac12aT^2.
\]
導入での$v(t)=2t$は,\,$v(0)=0$,\,$a=2$の場合なので,\,$D=T^2$が再び得られる.最初の長方形の和,速度のグラフの三角形,部分積分という三つの計算が,同じ位置の変化を求めている.

この式は,単に物理の公式を増やすためのものではない.積分する関数を$v$から$v'$へ替えたことで,初めの値と,途中で変わった分とに分けて表せたのである.加速度が与えられている問題では,速度を途中で求めずに使えることもある.

